\documentclass[11pt]{article}
\usepackage[margin=1in]{geometry}
\usepackage[T1]{fontenc}
\usepackage{amsmath,amssymb,amsthm,mathtools}
\usepackage[colorlinks=true,linkcolor=blue,urlcolor=blue,citecolor=blue]{hyperref}
\newtheorem{theorem}{Theorem}
\newtheorem{lemma}{Lemma}
\newcommand{\E}{\mathbb E}
\newcommand{\Prob}{\mathbb P}
\newcommand{\Var}{\operatorname{Var}}
\setlength{\emergencystretch}{2em}
\title{Real zeros of balanced arrays when the number of columns grows at least linearly}
\author{ProgramComputer}
\date{November 26, 2026}
\begin{document}
\maketitle

\begin{abstract}
Let $a_k$ count labeled binary $n\times d$ arrays with column sums
all equal to $k$ and row sums at most $r$. We prove that for every
$\epsilon>0$, all sufficiently large $n$, every $d\ge\epsilon n$, and every
$1\le r\le d$, the polynomial $\sum_k a_kq^k$ has
$\lfloor nr/d\rfloor$ simple negative zeros. The threshold is
existential. A uniform margin enumeration reduces the local coefficient
ratios to those of a scalar capped-binomial polynomial. Near saturation
the ratios give direct term separation. Elsewhere the alternating
evaluation tends to a positive Jacobi theta value.
\end{abstract}

\begin{theorem}\label{thm:main}
For every $\epsilon>0$ there is an integer $N(\epsilon)$ such that,
whenever $n\ge N(\epsilon)$, $d\ge\epsilon n$, and $1\le r\le d$ are
integers, the balanced-array polynomial $P(q)=\sum_k a_kq^k$ has exactly
$\lfloor nr/d\rfloor$ distinct, strictly negative zeros.
\end{theorem}

This proves a growing-dimension range of Alexandersson's balanced-array
conjecture for uniform matroids \cite{Alex}. The theorem concerns a
single row count and gives no assertion about interlacing between
consecutive counts. Its argument supplies no numerical value of
$N(\epsilon)$.

\section{Normalized log-concavity and test points}

Write $D=\lfloor nr/d\rfloor$. Capacity gives $a_k=0$ for $k>D$, and
every quota $0\le k\le D$ is realized by taking $dk$ consecutive
labels from a cyclic list of the $n$ rows and using successive groups
of $k$ labels as columns. Each group has distinct labels and each row
is used at most $\lceil dk/n\rceil$ times. In particular $a_0=1$.

We recall the repository's normalized log-concavity reduction
\cite{Orient}, based on an orientation inequality of Kahn and
Neiman \cite[Lemma~8]{KN}. It gives
\begin{equation}\label{eq:gamma}
\Gamma_k:=\frac{a_k^2}{a_{k-1}a_{k+1}}
\ \ge\ \left(1+\frac1k\right)^d
          \left(1+\frac1{n-k}\right)^d
\ \ge\ (1+2/n)^{2d}\qquad(1\le k<D).
\end{equation}
Here are the details needed for this reduction. For a quota vector
$\mathbf{k}$, let $A(\mathbf{k})$ be the capped count and
$p(\mathbf{k})=A(\mathbf{k})/\prod_j\binom n{k_j}$ the cap probability
under independent uniform columns. The orientation inequality states
that the number of orientations with vertexwise lower bounds
$(\mathbf a,\mathbf b)$ on out- and indegrees is at most the number with
bounds $(\mathbf c,\mathbf e)$ if
$\mathbf a\ge\mathbf c,\mathbf e$ and
$\mathbf a+\mathbf b\ge\mathbf c+\mathbf e$.

For a pair of arrays, color their entries red and blue. Fix the common
and exclusive rows in every column. If the quotas agree, each column
has $m_j$ red-only and $m_j$ blue-only rows. Pair those rows in that
column and orient every pair from red to blue. This gives a multigraph
on row labels. Each array pair has $\prod_j m_j!$ pairings, a factor
depending only on the fixed common and exclusive data. At a row with
common count $c$ and exclusive degree $H$, upper caps $U_R,U_B$ give
out- and indegree lower bounds $\max(0,H+c-U_B)$ and
$\max(0,H+c-U_R)$. Therefore when $U_R\ge U_B$ at each row, swapping
the caps at any subset of rows can only increase the orientation
count. This compares the original array counts after summing over the
fixed data.

For the own-coordinate comparison, fix a column $i$ and let $F(S)$
count the other columns with cap $r-1$ on $S$ and cap $r$ elsewhere.
The cap comparison gives
$F(S\cap T)F(S\cup T)\le F(S)F(T)$, taking red lowered on $S\cap T$
and blue lowered on $S\cup T$ and swapping on $S\setminus T$.
By row symmetry $F(S)$ depends only on $|S|$. Taking $S,T$ of size
$q+1$ with intersection of size $q$ gives the adjacent log-concavity
inequality after conditioning on the occupied set in column $i$.

For distinct $i,j$, compare
$A(\mathbf{k}) A(\mathbf{k}+\mathbf e_i+\mathbf e_j)$ with
$A(\mathbf{k}+\mathbf e_i)A(\mathbf{k}+\mathbf e_j)$.
Make the two copies' quotas agree by adding a distinct dummy row to
each of these two columns and filling its cell in the lower-quota copy
only. In the first product the red copy has both dummy entries. In the
second the dummy in column $i$ changes to blue while the one in column
$j$ stays red. Every real row has equal degree bounds and the changed
dummy swaps $(1,0)$ to $(0,1)$. The same orientation comparison proves
that the first product is at most the second. The products of the
normalizing binomials agree. Thus each extension ratio
$p(\mathbf{k}+\mathbf e_i)/p(\mathbf{k})$ is nonincreasing in every
coordinate wherever the ratios are positive. Comparing the successive
extensions from $(k-1,\ldots,k-1)$ to $(k,\ldots,k)$ with the same
extensions from $(k,\ldots,k)$ to $(k+1,\ldots,k+1)$ proves log-concavity
of $a_k/\binom nk^d$. Equation \eqref{eq:gamma} follows by taking the
binomial ratios and using $k(n-k)\le n^2/4$.

Set $t_k=\sqrt{a_{k-1}/a_{k+1}}$. The strict decrease of adjacent
coefficient ratios makes the internal $t_k$ strictly increasing.
At $-t_k$ both neighboring terms have magnitude
$a_kt_k^k/\sqrt{\Gamma_k}$, and moving farther out on either side
decreases the magnitudes. In the sign of the central term, each
alternating tail starting at distance two is nonnegative by pairing
successive terms. In particular $\Gamma_k>4$ implies
$(-1)^kP(-t_k)>0$. If that sign holds at every internal node, the
positive constant term and the sign $(-1)^D$ for sufficiently large
$t$ give $D$ disjoint sign-change intervals. The degree then forces
exactly $D$ simple negative roots. Degrees zero and one follow directly.

\section{Reduction to a scalar coefficient}\label{sec:counting}

If Theorem~\ref{thm:main} failed, there would be a sequence with
$n\to\infty$, $d/n\ge\epsilon$, and an internal node $k$ in each row
with $(-1)^kP(-t_k)\le0$. In particular $\Gamma_k\le4$.
Equation \eqref{eq:gamma} bounds $d/n$ from above.
After passing to a subsequence write
\[
d/n\to\delta\in(0,\infty),\qquad r/d\to\rho,\qquad x=k/n\to x_0.
\]
The bounds $\log(1+1/k)\ge 1/(k+1)$ and the corresponding bounds
at $n-k$ show that $0<x_0<1$. Capacity gives $x_0\le\rho$.
We work throughout at comparable dimensions and densities in a compact
subset of $(0,1)$. The next estimate holds at all feasible quotas in
this range, including the maximum.

Put $\mu=dx$, $H(z)=\sum_{s=0}^r\binom ds z^s$,
$A_k=[z^{dk}]H(z)^n$, and $K_k=\binom nk^d/\binom{nd}{dk}$.
For $0<\mu<r$, tilt a row size $S$ by weights $\binom ds z^s$,
using the unique $z>0$ with mean $\mu$. The mean has derivative
$\Var(S)>0$ with respect to $\log z$ and ranges from 0 to $r$.
At $\mu=r$, take $S=r$ deterministically. Set $V=\Var(S)$.

\begin{lemma}\label{lem:enum}
In the stated compact range,
\begin{equation}\label{eq:enum}
a_k=K_kA_k\exp\left(-\tfrac12+\frac{V}{2d x(1-x)}+o(1)\right).
\end{equation}
\end{lemma}
\begin{proof}
First choose uniformly from the actual arrays and condition on all
rows except a pair. In a column requiring two ones from that pair,
the two rows receive one each. If $b$ columns require just one one,
their allocation gives a row-size difference $2Y-b$, where $Y$ is
binomial $(b,1/2)$ conditioned to the central symmetric interval
imposed by the two caps. This conditioning can only decrease the
absolute gap in stochastic order. The binomial tail bound gives
$\Prob(|2Y-b|>L)\le2e^{-L^2/(2d)}$, with zero gap when $b=0$.
A row more than $L$ from the average $\mu$ requires a pair with gap
larger than $L$. Therefore the fraction of actual arrays having a row
deviation larger than $d^{1/2+\eta}$ is $o(1)$ for any fixed $\eta>0$.

For each remaining row profile $(s_i)_{i=1}^n$, the column margins
are all $k$ and the row deviations have sum zero. Theorem~1 of
Canfield, Greenhill and McKay \cite{CGM}, with its small deviation
exponent chosen to be $\eta$, gives the number of arrays with this
profile as
\begin{equation}\label{eq:cgm}
K_k \prod_{i=1}^n\binom d{s_i}
 \exp\left(-\tfrac12+\frac{W}{2x(1-x)}+o(1)\right),
\qquad
W=\frac1{nd}\sum_{i=1}^n(s_i-\mu)^2 .
\end{equation}
Indeed the margins are positive for large $n$, the column deviations
vanish, and the row deviations are at most $d^{1/2+\eta}$. The two
dimension-growth conditions in that theorem hold because $n\asymp d$
and $x(1-x)$ is bounded below. Its logarithmic density condition has
bounded left side for the same reasons. The error is uniform over
these profiles.

For independent tilted sizes $S_i$ conditioned on total $dk$, the
profile has probability $\prod_i\binom d{s_i}/A_k$. Let us justify
the conditioning and averaging in \eqref{eq:cgm}. First, a
log-concave probability mass function with interval support has
$\log\Prob(X=\E X)\ge-\operatorname{Ent}(X)$ at an integer mean.
Apply Jensen's inequality to the piecewise linear interpolation of
the log mass. Convolution preserves such log-concavity. One proof
notes that the Toeplitz matrix of the sequence, padded by zeros,
has nonnegative order-two minors: comparing two terms of a fixed
sum of indices reduces this to the decreasing successive ratios.
Matrix multiplication represents convolution, so Cauchy--Binet
proves the assertion. Thus the conditioning event has probability
at least $1/(nd+1)$, since the tilted sum has interval support in
$[0,nd]$ and entropy at most $\log(nd+1)$. It is certain in the
deterministic case.

Next, $S$ is dominated in convex order by a binomial with the same
mean. The ratio of their masses is geometric up to $r$ and zero
beyond it, so the set where the first mass is at least the second
is an interval. For a convex test function subtract the affine
secant between the two endpoints of that interval, chosen between
lattice sites. The remaining function is nonpositive inside and
nonnegative outside. Since the laws have equal mass and mean,
the secant cancels and the stated comparison follows. Identical
or deterministic laws satisfy the comparison directly.
Comparison of exponential moments and the binomial Hoeffding
bound now give $\Prob(|S-\mu|>y)\le2e^{-2y^2/d}$.

Consequently $Y_i=(S_i-\mu)^2/d$ have a bounded exponential
moment at a fixed positive argument. Here $W=n^{-1}\sum_iY_i$
has mean $V/d$. For every fixed $u>0$, the centered exponential
moment has logarithm $O(t^2)$ on a fixed neighborhood of zero,
uniformly in the parameters. Exponential Markov then gives
$\Prob(|W-V/d|>u)\le 2e^{-c_u n}$. For each fixed $C>0$,
$\E e^{CW}=(\E e^{CY_1/n})^n$ is bounded. Also the probability
of an irregular row is at most $2n e^{-2d^{2\eta}}$.
Cauchy--Schwarz shows that either exceptional event contributes
$o(1)$ to the exponentially weighted conditional expectation,
even after division by the conditioning probability. On the
complement $W$ is within $u$ of $V/d$. Letting $u$ decrease to zero
gives conditional average $\exp(V/[2d x(1-x)])+o(1)$ over the
regular profiles. Summing \eqref{eq:cgm} and using the relative
tail bound in the actual counting measure proves \eqref{eq:enum}.
\end{proof}

At any bounded integer shift of quota from $k$, the centered
log-curvature $2\log K_k-\log K_{k-1}-\log K_{k+1}$ tends to zero.
For the numerator $\binom nk^d$ its value tends to
$\delta/[x_0(1-x_0)]$ by the adjacent-binomial ratios. The curvature
ratio of the denominator $\binom{nd}{dk}$ is
\[
\prod_{j=1}^d
\left(1+\frac{d}{dk-d+j}\right)
\left(1+\frac{d}{nd-dk-d+j}\right).
\]
Taking logarithms and summing gives the same limit.

\section{Vanishing scaled slack}\label{sec:endpoint}

Put $\lambda=d/n$ and $h=r-\mu$. Consider $h=o(\sqrt d)$, so that
$r/d\to x_0\in(0,1)$. We will show
\begin{equation}\label{eq:edge-curv}
\log\frac{A_k^2}{A_{k-1}A_{k+1}}\longrightarrow+\infty,
\qquad V/d\longrightarrow0
\end{equation}
at the three quotas in the curvature comparison. The statement about
curvature is for an internal $k$, so $h-\lambda\ge0$.

Set $a=r/(d-r+1)$ and $H(z)=\binom dr z^r B(a/z)$, where
\[
B(q)=\sum_{j=0}^r c_jq^j,\qquad
c_j=\prod_{i=0}^{j-1}\frac{1-i/r}{1+i/(d-r+1)}.
\]
Extend the coefficients by zero. They decrease, $c_0=c_1=1$, and
$0\le1-c_j\le Cj(j-1)/d$: for $j\le r$, sum the deficits of the
factors, and for $j>r$ enlarge $C$ if necessary. A tilted deficit
$J=r-S$ has mass proportional to $c_jq^j$. Compare it for $q<1$
with a geometric law of mean $w=q/(1-q)$. Under that law the
normalization of the density factor $c_J$ is $1+O(w^2/d)$.
For $p=1,2$ and $w=o(\sqrt d)$, the tilted moment differs from the
geometric moment by $O(w(1+w)^{p+1}/d)$. This follows from the
coefficient bound and the geometric factorial moments
$\E(J)_b=b!w^b$. Also the tilted mean is at most $w$ by the
nonpositive covariance of the increasing function $J$ and the
decreasing function $c_J$. Trial geometric means $h,2h$ therefore
bracket the desired tilt $q_h$ for positive slack, uniformly in the
stated range. Applying the same moment comparisons at that tilt gives
\begin{equation}\label{eq:edge-var}
w=h\bigl(1+O((1+h)^2/d)\bigr),\qquad
V=h(1+h)\bigl(1+O((1+h)^2/d)\bigr).
\end{equation}
In the variance estimate use the geometric variance $w(1+w)$ and
subtract the square of the expectation. At zero slack the tilt is
$q_h=0$ and $V=0$.

We need a central probability estimate that works even for bounded
total deficit. For a log-concave law on consecutive integers with
variance $\sigma^2$ and integer mean, the mass at its mean is
comparable to $(1+\sigma)^{-1}$ with numerical constants. For the
lower bound, the entropy inequality used in Lemma~\ref{lem:enum}
applies. Adding an independent uniform variable on $[-1/2,1/2]$
preserves entropy, now in its differential version, and gives variance
$\sigma^2+1/12$. The Gaussian upper bound for differential entropy
follows by nonnegativity of relative entropy to the Gaussian of the
same mean and variance. For the upper bound let $b$ be the modal
mass. On each side of the mode the first distance $L$ with mass at
most $b/2$ satisfies $L\le 2/b+1\le3/b$, extending by zero off the
support. Log-concavity bounds the mass at distance $j\ge L$ by
$b\,2^{-j/L}$. At smaller distances use the modal bound. Taken
together the masses on each side are bounded by
$2b\, 2^{-bj/3}$. Summing $j^2$ times this majorant gives
$\sigma^2\le C/b^2$, and $b\le1$ yields the upper bound.

Write $\ell=nh$ and $G(h)=\log B(q_h)-h\log q_h$, with $G(0)=0$.
The exact coefficient expression is
\[
A_k=\binom dr^{\,n}a^\ell e^{nG(h)}
       \Prob(J_1+\cdots+J_n=\ell).
\]
The deficit law is log-concave by its successive ratios, as is its
convolution. The central probability estimate and \eqref{eq:edge-var}
give
\begin{equation}\label{eq:edge-log}
\log A_k=n\log\binom dr+\ell\log a+nG(h)
   -\tfrac12\log(1+n h(1+h))+O(1).
\end{equation}
At zero slack the probability is one and $A_k=\binom dr^n$
exactly. This covers the possibility that the maximum feasible quota
is saturated, without presuming a diverging total variance.

For $h>0$, differentiation gives $G'(h)=-\log q_h$ and
$G''(h)=-1/V(h)$. The entropy is continuous at zero because
$q_h\sim h$ there. When $h\to\infty$, the logarithmic prefactors in
\eqref{eq:edge-log} contribute $O(1)$ to curvature since the
neighboring variance scales are comparable. Meanwhile
\[
2G(h)-G(h-\lambda)-G(h+\lambda)
  =\int_{-\lambda}^{\lambda}\frac{\lambda-|u|}{V(h+u)}\,du
  \sim \frac{\lambda^2}{h(1+h)}.
\]
Multiplication by $n$ gives a diverging curvature. For bounded $h$,
the integral is bounded below by a positive constant on its central
subinterval $|u|\le\lambda/2$, by \eqref{eq:edge-var} and
$\lambda\asymp1$. The bound extends to $h-\lambda=0$ by continuity
of $G$, taking a limit from positive slack. The remaining terms in
\eqref{eq:edge-log} contribute $O(\log n)$ to curvature. This proves
\eqref{eq:edge-curv}, including bounded or vanishing slack at the
higher adjacent quota.

\section{Positive scaled slack}\label{sec:central}

Assume now that $(r-\mu)/\sqrt d\to H_*\in(0,\infty]$.
We show that $V/d\to v>0$ and that the tilted central mass is
asymptotic to $(2\pi n d v)^{-1/2}$. The same variance limit holds
under any bounded real quota shift.

Write the row tilt as $z=e^\eta$. The capped law is a
binomial of parameter $e^\eta/(1+e^\eta)$, conditioned at the cap.
At the trial values $\eta=\log(x/(1-x))+\tau/\sqrt d$, with $\tau$
bounded, the standardized unrestricted row size $(S-\mu)/\sqrt d$
tends to a normal of mean $\tau v_0$ and variance $v_0$ when
$\tau$ converges, where $v_0=x_0(1-x_0)$. The binomial central
limit theorem and Hoeffding bounds give convergence of the
capped variable and all its fixed moments to that normal
conditioned to be at most $H_*$. The cap probability is bounded
away from zero for $\tau$ in compact sets, by the positive
limiting normal left tail and a subsequence argument. If
$H_*=\infty$, the conditioning disappears.

For finite $H_*>0$, the limiting capped law has density proportional
to $\exp(\tau y-y^2/(2v_0))$ on $y\le H_*$. Its mean increases
strictly and continuously from $-\infty$ to $H_*$ as $\tau$ increases.
For $\tau\to-\infty$ center at the unrestricted mean, where truncation
is negligible. For $\tau\to+\infty$ the distance $H_*-y$ has density
proportional to
$\exp(-(\tau-H_*/v_0)(H_*-y)-(H_*-y)^2/(2v_0))$, and rescaling by
$\tau-H_*/v_0$ shows its expectation tends to zero. In the infinite
case the normal mean increases through zero at $\tau=0$.
In either case choose two finite trials with opposite mean signs.
They bracket the actual tilt for large $n$, and its $\tau$ tends to
the unique zero of the limiting mean. The limiting variance $v$
is strictly positive. Bounded shifts of quota alter the scaled
cap position by $O(\sqrt d/n)=o(1)$ and preserve $x_0$, so the
same sequential argument gives the asserted uniformity in such
shifts. The number of available quota steps above and below
$k$ tends to infinity, using
$D-k=\lfloor n(r-\mu)/d\rfloor$ above and $k/n\to x_0>0$ below.

For the local estimate let $\phi(t)=\E e^{it(S-\mu)}$. At an
integer total $n\mu=dk$, Fourier inversion gives the central
probability as $(2\pi)^{-1}\int_{-\pi}^{\pi}\phi(t)^n\,dt$.
With an independent copy $S'$, the identity
$1-|\phi(t)|^2=\E(1-\cos(t(S-S')))$ and the variance lower and
fourth moment upper bounds give $|\phi(t)|\le e^{-c d t^2}$
for $|t|\le a_0/\sqrt d$ and some fixed $a_0>0$. Taylor expansion
using the bounded standardized third moment gives
$\phi(u/\sqrt{nd})^n\to e^{-v u^2/2}$ for fixed $u$.
On $a_0\le\sqrt d|t|\le M$ the limiting centered continuous
law has characteristic function of modulus strictly less than
one. Convergence is uniform on this compact frequency interval
by the uniform first-moment bound. On $M/\sqrt d\le |t|\le\pi$,
summation by parts gives
\[
|\phi(t)|\le \frac{2\max_s\Prob(S=s)}{|1-e^{it}|}
            \le\frac{C}{\sqrt d\,|t|}.
\]
Here unimodality of the capped binomial gives total variation
twice the maximum mass after padding with zeros. Stirling's bounds
and the cap probability bound give maximum mass $O(d^{-1/2})$.
Take $M$ large. Both outer intervals are negligible after multiplication
by $\sqrt{nd}$, and Gaussian domination on the first interval gives
\begin{equation}\label{eq:lclt}
\Prob(S_1+\cdots+S_n=dk)\sim(2\pi nd v)^{-1/2}.
\end{equation}
The same equivalent holds at each fixed integer quota shift.

\section{The limiting alternating sign}

Return to the failed nodes. Pass to a subsequence on which the
scaled slack converges in $[0,\infty]$. If its limit is zero,
\eqref{eq:edge-curv}, \eqref{eq:enum} and the cancellation of
the $K_k$ curvature give $\Gamma_k\to\infty$, contradicting
$\Gamma_k\le4$.

For a positive or infinite limit, define
$E(k)=n\log H(z)-dk\log z$ at the row tilt, allowing real quota
shifts locally. Differentiation gives $E''(k)=-d^2/(nV)$,
which tends uniformly on bounded shifts to $-\delta/v$.
For integers $k$, the tilted representation is
$A_k=e^{E(k)}\Prob(\sum_i S_i=dk)$. Equation \eqref{eq:lclt}
gives a common asymptotic prefactor at bounded shifts. The
variance correction in \eqref{eq:enum} also has a common limit,
and the $K_k$ curvature tends to zero. Integrating $E''$ against
the triangular second-difference kernel therefore gives, for
each fixed integer $j$,
\[
\log\Gamma_{k+j}\longrightarrow g=\delta/v>0.
\]

At $-t_k$, divide the magnitude of every term by $a_kt_k^k$.
The magnitudes at distance $j\ge1$ on the right and left are,
respectively,
\[
\Gamma_k^{-j/2}\prod_{i=1}^{j-1}\Gamma_{k+i}^{-(j-i)},
\qquad
\Gamma_k^{-j/2}\prod_{i=1}^{j-1}\Gamma_{k-i}^{-(j-i)},
\]
with absent terms assigned zero. These identities follow by
multiplying the successive adjacent ratios. For fixed $j$ they
converge to $e^{-g j^2/2}$. Equation \eqref{eq:gamma} also gives
$\log\Gamma_i\ge c_0>0$ uniformly in all internal $i$ along the
sequence for large $n$. The magnitudes are thus dominated by
$e^{-c_0j^2/2}$, a summable bound. Dominated convergence gives
\[
\lim\frac{(-1)^kP(-t_k)}{a_kt_k^k}
=\sum_{j\in\mathbb Z}(-1)^j e^{-g j^2/2}
=\prod_{m=1}^{\infty}(1-e^{-gm})
                        (1-e^{-g(m-1/2)})^2>0.
\]
The identity is the Jacobi triple product. Its factors are positive
and their deficits are summable. Equivalently, Poisson summation
expresses the value as
$\sqrt{2\pi/g}\sum_{m\in\mathbb Z}e^{-2\pi^2(m+1/2)^2/g}$.
This contradicts the failure of the node in the other possible
regime as well. The sign test and the degree argument finish the
proof of Theorem~\ref{thm:main}.

\begin{thebibliography}{9}
\raggedright
\bibitem{Alex}
P.~Alexandersson,
\emph{Balanced arrays of uniform matroids},
in \emph{Real-rooted polynomials: open problems},
The Symmetric Functions Catalog, 2026.
\href{https://www.symmetricfunctions.com/realRootedOpenProblems.htm\#conjBalancedArrays}
{Original conjecture}.

\bibitem{CGM}
E.~R. Canfield, C.~Greenhill, and B.~D. McKay,
Asymptotic enumeration of dense $0$--$1$ matrices with specified line sums,
\emph{J. Combin. Theory Ser. A} \textbf{115} (2008), 32--66.
\href{https://doi.org/10.1016/j.jcta.2007.03.009}
{doi:10.1016/j.jcta.2007.03.009}.

\bibitem{KN}
J.~Kahn and M.~Neiman,
\emph{Conditional negative association for competing urns},
\emph{Random Structures Algorithms} \textbf{41} (2012), 262--281.
\href{https://doi.org/10.1002/rsa.20395}{doi:10.1002/rsa.20395}.
The orientation inequality is Lemma~8 in
\href{https://arxiv.org/abs/1001.0610}{arXiv:1001.0610}.

\bibitem{Orient}
ProgramComputer,
\emph{An orientation proof of normalized log-concavity for balanced arrays},
Open Math Research, 2026.
\href{https://github.com/ProgramComputer/open-math-research/blob/research/topics/balanced-arrays/orientation-normalization.md}
{Repository proof}.
\end{thebibliography}
\end{document}
